\section{Waiting and Signalling Using JVM Monitors}

We now consider the second aspect of JVM monitors, namely the wait signal
mechanism.  If |o| is a reference object, then it has the following
operations, each of which can be performed only inside a |synchronized| block
on~|o|:
\begin{itemize}
\item |o.wait()|: suspend, giving up the lock on~|o|; wait for a signal, and
  then re-acquire the lock;

\item |o.notify()| : signal to a thread waiting on~|o|, if there is
  one (if there is no thread waiting, the operation has no effect);

\item |o.notifyAll()|: signal to all threads waiting on~|o|. 
\end{itemize}
%
Each operation should be performed while holding the lock on~|o|, i.e.~from
inside a |synchronized| block on~|o|.  When a thread signals, it keeps the
lock: the thread that receives the signal must wait for the signaller to
release the lock, and then needs to compete with other threads to re-acquire
the lock.

The operation |wait()| is shorthand for |this.wait()|, i.e.~where the thread
waits on the current object.  Similar things hold for |notify()| and
|notifyAll()|. 

Figure~\ref{fig:partial-queue-JVM-monitor} illustrates the mechanism in an
implementation of a partial queue.  The |dequeue| operation has to wait if the
queue is empty.  An |enqueue| operation performs a |notify|, to wake up a
waiting |dequeue|.

%%%%%

\begin{figure}
\begin{scala}
import scala.collection.mutable.Queue

class JVMMonitorPartialQueue[T] extends PartialQueue[T]{
  /** The current state of the queue. */
  private val queue = new Queue[T]

  /** Enqueue £x£. */
  def enqueue(x: T) = synchronized{ 
    queue.enqueue(x) 
    notify() // Signal to a waiting £dequeue£.
  }

  /** Try to dequeue, blocking while the queue is empty.  */
  def dequeue(): T = synchronized{ 
    while(queue.isEmpty) wait() // Wait for a value.
    queue.dequeue() 
  }
}
\end{scala}
\caption{A partial queue implemented using a JVM monitor.}
\label{fig:partial-queue-JVM-monitor}
\end{figure}

%%%%%

As with the implementation using a |Condition| in
Figure~\ref{fig:MonitorPartialQueue}, a waiting |dequeue| needs to recheck
whether the queue is empty when it receives a signal, because it is possible
that some other thread has performed a |dequeue| in the meantime and emptied
the queue.  

However, there is another reason why it needs to perform this check, which
would hold even if there were no other thread that could attempt a |dequeue|.
Sometimes a call of |wait| wakes up even if no other thread has signalled!
This is known as a \emph{spurious wakeup}.  I consider this a bug (apparently
it is caused by a similar feature in operating system threading
implementations; but it is a feature that the JVM fails to deal with).
Nevertheless, we have to deal with it.  That means that whenever a call of
|wait| wakes up, it should recheck the relevant condition, in case it has
received a spurious wakeup.

\begin{instruction}
Study the details of the implementation.
\end{instruction}

The implementation of a bounded partial queue is fairly straightforward,
because there is a single point in the code where a thread can wait.  However,
things are harder when threads can wait in multiple different places.  SCL
monitors allow signals to be targeted to the correct place, by using a
different |Condition| for each.  However, there is no similar way to target
signals in a JVM monitor.  If there are threads waiting in multiple places,
when there is a signal, it will be nondeterministic where the signal is
received, so it might well be received in the wrong place!  

In such circumstances it is necessary to use |notifyAll|.  This means that the
signal will reach the right place, if there is a thread waiting there.
However, threads waiting in other places should identify that the signal was
not intended for them, and wait again.  And likewise, if there are multiple
threads waiting at the place to which the signal was intended, the first to
obtain the lock can continue (if not preempted by another thread), but the
remainder should normally wait again.  

This is illustrated in Figure~\ref{fig:bounded-partial-queue-JVM-monitor}, to
implement a bounded partial queue.  Again, a |dequeue| has to wait when the
queue is empty; in addition, an |enqueue| has to wait when the queue is full.
An |enqueue| signals to a waiting |dequeue|; but it has to perform
|notifyAll()|, in case there are both |enqueue|s and |dequeue|s waiting.
Likewise, a |dequeue| signals to a waiting |enqueue|, but again has to perform
|notifyAll()|.  (In fact, when we use |notifyAll|, we can't have enqueues and
dequeues waiting simultaneously, because the waiting conditions are disjoint;
but we could if we used |notify|.)

%%%%%

\begin{figure}
\begin{scala}
class JVMMonitorBoundedPartialQueue[T](bound: Int) extends PartialQueue[T]{
  require(bound > 0)

  /** The current state of the queue. */
  private val queue = new Queue[T]

  /** Enqueue £x£. */
  def enqueue(x: T) = synchronized{ 
    while(queue.length == bound) wait() // Wait for a space (1).
    queue.enqueue(x)
    notifyAll() // Signal to a waiting dequeue at (2).
  }

  /** Try to dequeue, blocking while the queue is empty.  */
  def dequeue(): T = synchronized{ 
    while(queue.isEmpty) wait() // Wait for a value (2).
    val result = queue.dequeue() 
    notifyAll() // Signal to a waiting enqueue at (1).
    result
  }
}
\end{scala}
\caption{A bounded partial queue implemented using a JVM monitor.}
\label{fig:bounded-partial-queue-JVM-monitor}
\end{figure}

%%%%%

\begin{instruction}
Study the implementation.  Make sure you understand why using |notify| could
go wrong.
\end{instruction}

%% We can slightly improve the efficiency of the implementations.  We can arrange
%% for an |enqueue| to perform a |notifyAll| only when it makes the queue
%% non-empty, i.e.~if it was previously empty:
%% %
%% \begin{scala}
%%   def enqueue(x: T) = synchronized{ 
%%     while(queue.length == bound) wait() // Wait for a space (1).
%%     queue.enqueue(x)
%%     if(queue.length == 1) notifyAll() // Signal to a waiting dequeue at (2).
%%   }
%% \end{scala}
%% %
%% This works because we have an invariant that there is a waiting |dequeue| only
%% if the queue is empty; thus there is no point in signalling at other times. 
%% % Actually, this doesn't save much -- a notifyAll when there are no waiting
%% % threads. 
%% \begin{scala}
%%   def dequeue(): T = synchronized{ 
%%     while(queue.isEmpty) wait() // Wait for a value (2).
%%     val result = queue.dequeue() 
%%     if(queue.length == bound-1) 
%%       notifyAll() // Signal to a waiting enqueue at (1).
%%     result
%%   }
%% \end{scala}

Having to use |notifyAll| is a disadvantage of JVM monitors.  It will
typically mean that several threads are woken up, only for them to have to
wait again.  In particular, scheduling and descheduling threads are quite
expensive operations, to this adds a fairly significant overhead.  In such
cases, it is probably better to use |Condition|s, which allow targetted
signalling. 

%%%%%%%%%%%%%%%%%%%%%%%%%%%%%%%%%%%%%%%%%%%%%%%%%%%%%%%
